\inicializarSeccion{Campo eléctrico generado}

\tab A continuación se muestra el campo eléctrico generado por la distribución de cargas positivas y negativas descrita en las secciones anteriores

Iniciaremos con las simulaciones que contienen las mismas cargas (8 cargas puntuales en horizontal) pero con diferentes distancias entre las líneas de cargas.

\begin{figure}[H]
    \centering
    \begin{minipage}[t]{0.48\textwidth}
        \centering
        \includesvg[width=1\textwidth, inkscapelatex=false]{\nat{parte2/casoOchoCinco.svg}}
        \caption{Simulación del campo eléctrico generado por un dipolo con separación de $d=0.05$ con 8 cargas puntuales horizontales.}
        \label{fig:casoOchoCinco}
    \end{minipage}\hfill
    \begin{minipage}[t]{0.48\textwidth}
        \centering
        \includesvg[width=1\textwidth, inkscapelatex=false]{\nat{parte2/casoOchoDiez.svg}}
        \caption{Simulación del campo eléctrico generado por un dipolo con separación de $d=0.10$ con 8 cargas puntuales horizontales.}
        \label{fig:casoOchoDiez}
    \end{minipage}

    \vspace{0.5cm}

    \begin{minipage}[t]{0.48\textwidth}
        \centering
        \includesvg[width=1\textwidth, inkscapelatex=false]{\nat{parte2/casoOchoVeinte.svg}}
        \caption{Simulación del campo eléctrico generado por un dipolo con separación de $d=0.20$ con 8 cargas puntuales horizontales.}
        \label{fig:casoOchoVeinte}
    \end{minipage}
\end{figure}

Otros casos que podemos observar es con una cantidad distinta de cargas:

\begin{figure}[H]
    \centering
    \begin{minipage}[t]{0.48\textwidth}
        \centering
        \includesvg[width=1\textwidth, inkscapelatex=false]{\nat{parte2/casoTresCinco.svg}}
        \caption{Simulación del campo eléctrico generado por un dipolo con separación de $d=0.05$ con 3 cargas puntuales horizontales.}
        \label{fig:casoTresCinco}
    \end{minipage}\hfill
    \begin{minipage}[t]{0.48\textwidth}
        \centering
        \includesvg[width=1\textwidth, inkscapelatex=false]{\nat{parte2/casoVeinteCinco.svg}}
        \caption{Simulación del campo eléctrico generado por un dipolo con separación de $d=0.05$ con 20 cargas puntuales horizontales.}
        \label{fig:casoVeinteCinco}
    \end{minipage}
\end{figure}

Por ultimo, un caso mas en particular es si hay un número distinto de cargas positivas y negativas: 

\begin{figure}[H]
    \centering
    \begin{minipage}[t]{0.48\textwidth}
        \centering
        \includesvg[width=1\textwidth, inkscapelatex=false]{\nat{parte2/casoVeinte-TresQuince.svg}}
        \caption{Simulación del campo eléctrico generado por un dipolo con separación de $d=0.15$ con 20 cargas puntuales horizontales y 3 cargas puntuales horizontales.}
        \label{fig:casoVeinteTresQuince}
    \end{minipage}\hfill
    \begin{minipage}[t]{0.48\textwidth}
        \centering
        \includesvg[width=1\textwidth, inkscapelatex=false]{\nat{parte2/casoDos-DiezQuince.svg}}
        \caption{Simulación del campo eléctrico generado por un dipolo con separación de $d=0.15$ con 2 cargas puntuales horizontales y 10 cargas puntuales horizontales.}
        \label{fig:casoDosDiezQuince}
    \end{minipage}
\end{figure}